\section{Worked example}\label{app:worked-example}

At $m=5$, the admissible language $X_5$ has $F_7=13$ words and $\Fold_5$ groups the $32$ raw words by Fibonacci valuation modulo $13$.

With the convention
\[
\Val_5(\omega)=\sum_{k=1}^5 \omega_k F_{k+1},
\]
the raw word $\omega=11010$ has valuation
\[
\Val_5(11010)=F_2+F_3+F_5=1+2+5=8.
\]
The unique admissible word of length five with the same residue is the word $00001$, for which $\Val_5(00001)=F_6=8$. Hence
\[
\Fold_5(11010)=00001.
\]

The whole fiber over $00001$ is
\[
\mathcal F_5(00001)=\{00001,\ 00110,\ 11010\}.
\]
All three words have Fibonacci valuation congruent to $8$ modulo $13$, and no other five-bit word does. This illustrates the two ingredients used in the main text: isolated windows can collapse onto the same admissible type, but the sequence-level ambiguity disappears once overlapping windows are read together.
